\oldpage[0\textsubscript{IV-1}]{75}

  \item[b)] L'homomorphisme
  $u_0:M/\mathfrak QM\to N/\mathfrak QN$ déduit de $u$ par passage aux
  quotients est inversible à gauche.
\end{enumerate}
\end{proposition}

\begin{proof}
On a vu (19.1.7) que la condition a) équivaut à dire que $u_\lambda$ est
inversible à gauche pour tout $\lambda$~; comme $\mathfrak Q$ est un idéal
ouvert, donc contient un $\mathfrak J_\lambda$, on en déduit aussitôt que
$u_0$ est inversible à gauche. Pour montrer inversement que b) entraîne a),
notons que pour tout $\lambda$, $\mathfrak Q/\mathfrak J_\lambda$ est par
hypothèse un idéal nilpotent de $A/\mathfrak J_\lambda$. Notre assertion
résultera de la proposition suivante.
\end{proof}

\begin{proposition}[19.1.10]
\label{0.19.1.10-fr}
Soient $A$ un anneau, $M$, $N$ deux $A$-modules, $N$ étant projectif,
$u:M\to N$ un $A$-homomorphisme. Soit $\mathfrak J$ un idéal de $A$~; on
suppose vérifiée l'une des conditions suivantes~:
\begin{enumerate}
  \item[(i)] $\mathfrak J$ est nilpotent.
  \item[(ii)] $\mathfrak J$ est contenu dans le radical de $A$ et $M$ est de
  type fini.
\end{enumerate}
Alors, pour que $u$ soit inversible à gauche, il faut et il suffit que
l'homomorphisme
$u_0:M/\mathfrak JM\to N/\mathfrak JN$ de $(A/\mathfrak J)$-modules, déduit
de $u$ par passage aux quotients, soit inversible à gauche.
\end{proposition}

\begin{proof}
La condition étant évidemment nécessaire, prouvons qu'elle est suffisante.
Soit $v_0$ un inverse à gauche de $u_0$~; l'homomorphisme composé
$N\to N/\mathfrak JN\xrightarrow{v_0}M/\mathfrak JM$ se factorise en
$N\xrightarrow{v}M\to M/\mathfrak JM$ puisque $N$ est projectif~; alors
$w=v\circ u$ est un endomorphisme de $M$ tel que l'endomorphisme $w_0$ de
$M/\mathfrak JM$ déduit par passage aux quotients soit l'identité, et il suffit
de prouver que $w$ est lui-même bijectif (car alors $w^{-1}\circ v$ sera un
inverse à gauche de $u$). Distinguons maintenant les deux cas.

(i) Pour tout $n$ on a le diagramme commutatif
\[
  \xymatrix{
    (\mathfrak J^n/\mathfrak J^{n+1})\otimes_{A/\mathfrak J}
      (M/\mathfrak JM)\ar[r]\ar[d]_{1\otimes\operatorname{gr}^0(w)} &
      \mathfrak J^nM/\mathfrak J^{n+1}M\ar[d]^{\operatorname{gr}^n(w)}\\
    (\mathfrak J^n/\mathfrak J^{n+1})\otimes_{A/\mathfrak J}
      (M/\mathfrak JM)\ar[r] & \mathfrak J^nM/\mathfrak J^{n+1}M
  }
\]
où les flèches horizontales sont \emph{surjectives}, et comme
$\operatorname{gr}^0(w)=w_0$ est l'identité, il en est de même de
$\operatorname{gr}^n(w)$ qui, \emph{a fortiori}, est bijectif. La filtration
$\mathfrak J$-préadique sur $M$ étant finie puisque $\mathfrak J$ est
nilpotent, on en conclut que $w$ est bijectif (Bourbaki, \emph{Alg. comm.},
chap.~III, §~2, n\textsuperscript{o}~8, cor.~3 du th.~1).

(ii) Il suffit de montrer que pour tout idéal maximal $\mathfrak m$ de $A$,
l'endomorphisme $w_{\mathfrak m}$ de $M_{\mathfrak m}$ est bijectif
(Bourbaki, \emph{Alg. comm.}, chap.~II, §~3, n\textsuperscript{o}~3, th.~1)
et comme $\mathfrak JA_{\mathfrak m}\subset\mathfrak mA_{\mathfrak m}$, et
$A_{\mathfrak m}/\mathfrak JA_{\mathfrak m}=(A/\mathfrak J)_{\mathfrak m}$,
on est ramené à prouver la proposition lorsque $A$ est un anneau local. En
outre, on peut supposer que $\mathfrak J$ est l'idéal maximal
$\mathfrak J_0$ de $A$, car si $u_0$ est inversible
\oldpage[0\textsubscript{IV-1}]{76}
à gauche, il en est de même de
$u_{00}:M/\mathfrak J_0M\to N/\mathfrak J_0M$ obtenu par tensorisation de
$u_0$ avec $1_{A/\mathfrak J_0}$, puisque l'on a
$M/\mathfrak J_0M=(M/\mathfrak JM)\otimes_{A/\mathfrak J}(A/\mathfrak J_0)$.
Supposons donc $\mathfrak J$ maximal, de sorte que $A/\mathfrak J$ est un
corps. Il suffit évidemment de montrer que $M$ est un $A$-module \emph{libre}
sous les conditions de l'énoncé~: en effet, $\det(w_0)$ est alors l'image
canonique de $\det(w)$ dans $A/\mathfrak J$, donc $\det(w)$ n'appartient pas à
l'idéal $\mathfrak J$ et est par suite inversible. Or le
$(A/\mathfrak J)$-espace vectoriel $M/\mathfrak JM$ étant libre de type fini,
il y a un $A$-module de type fini $L$ et un $A$-homomorphisme $f:L\to M$ tel
que l'homomorphisme $f_0:L/\mathfrak JL\to M/\mathfrak JM$ déduit de $f$ par
passage aux quotients soit bijectif. Comme $M$ est de type fini, on en conclut
tout d'abord que $f$ est \emph{surjectif} par le lemme de Nakayama (Bourbaki,
\emph{Alg. comm.}, chap.~II, §~3, n\textsuperscript{o}~2, cor.~1 de la
prop.~4)~; en outre, si $g=u\circ f$, l'homomorphisme
$g_0:L/\mathfrak JL\to N/\mathfrak JN$ déduit par passage aux quotients est
inversible à gauche, et puisqu'ici $L$ est libre, la remarque du début prouve
que $g$ est lui-même inversible à gauche~; or cela entraîne évidemment que
$f$ est \emph{injectif}, ce qui achève la démonstration.
\end{proof}

Mentionnons en passant les corollaires suivants de (19.1.10)~:

\begin{corollary}[19.1.11]
\label{0.19.1.11-fr}
Soient $A$ un anneau local, $k$ son corps résiduel, $M$ un $A$-module de type
fini, $N$ un $A$-module projectif, $u:M\to N$ un homomorphisme. Pour que $u$
soit inversible à gauche, il faut et il suffit qu'il existe un système de
générateurs $(x_i)_{1\leq i\leq m}$ de $M$ tel que les images par
$u\otimes 1:M\otimes_Ak\to N\otimes_Ak$ des $x_i\otimes1$ soient linéairement
indépendantes dans $N\otimes_Ak$~; les $x_i$ forment alors une base de $M$.
\end{corollary}

La condition est évidemment nécessaire, car si $u$ est inversible à gauche,
$M$ est un $A$-module projectif de type fini, donc libre. Inversement, si la
condition est satisfaite, il est clair que les $x_i\otimes1$ forment une base
de $M\otimes_Ak$ et que $u\otimes1$ est inversible à gauche~; il suffit donc
d'appliquer (19.1.10) à l'idéal maximal $\mathfrak J$ de $A$.

\begin{corollary}[19.1.12]
\label{0.19.1.12-fr}
Soient $A$ un anneau, $M$ un $A$-module de type fini, $N$ un $A$-module
projectif, $u:M\to N$ un homomorphisme. Pour tout idéal premier
$\mathfrak p\in\operatorname{Spec}(A)$, les conditions suivantes sont
équivalentes~:
\begin{enumerate}
  \item[\textup{a)}] L'homomorphisme
  $u_{\mathfrak p}:M_{\mathfrak p}\to N_{\mathfrak p}$ est inversible à
  gauche.
  \item[\textup{b)}] L'homomorphisme
  $u\otimes1:M\otimes_A\boldsymbol{k}(\mathfrak p)\to N\otimes_A\boldsymbol{k}(\mathfrak p)$ est
  injectif \textup{(on rappelle que $\boldsymbol{k}(\mathfrak p)$ désigne le corps résiduel de $A$
  en $\mathfrak p$)}.
  \item[\textup{c)}] Il existe un système fini d'éléments $x_i\in M$
  $(1\leq i\leq m)$ tel que les images des $x_i$ dans $M_{\mathfrak p}$
  engendrent $M_{\mathfrak p}$, et un système de $m$ formes linéaires
  $y_i^*$ sur $N$ $(1\leq i\leq m)$ tel que
  $\det(\mathopen{<}y_i^*,u(x_j)\mathclose{>})\notin\mathfrak p$.
  \item[\textup{d)}] Il existe $f\in A\setmin\mathfrak p$ tel que l'homomorphisme
  $u_f:M_f\to N_f$ soit inversible à gauche.
\end{enumerate}

En outre, l'ensemble des $\mathfrak p\in\operatorname{Spec}(A)$ vérifiant ces
conditions est ouvert dans $\operatorname{Spec}(A)$.
\end{corollary}

La dernière condition est conséquence triviale de \emph{d)}. Comme $N$ est
projectif, il est facteur direct d'un $A$-module libre $A^{(I)}$~; en outre,
comme $M$ est de type fini, $u(M)$ est contenu dans un sous-module de $A^{(I)}$
de la forme $A^n$ pour $n$ fini~; comme chacun des énoncés \emph{a)}, \emph{b)}, \emph{c)}, \emph{d)} est
équivalent à l'énoncé correspondant où l'on remplace $N$ par $N\oplus P$ ($u$
étant toujours supposé appliquer $M$ dans $N$), on est ramené au cas où $N$
est \emph{libre de type fini}. Il est trivial que \emph{d)} entraîne \emph{a)}, et \emph{a)} et \emph{b)}
sont équivalentes en vertu de (19.1.11)~; d'ailleurs \emph{a)} entraîne que
$M_{\mathfrak p}$ est libre (19.1.11), donc \emph{a)} entraîne \emph{c)}, en prenant pour les
$x_i$ une famille dont les images dans $M_{\mathfrak p}$ forment une base de ce
\oldpage[0\textsubscript{IV-1}]{77}
$A_{\mathfrak p}$-module et en notant que (puisque $N$ est libre de type fini)
toute forme linéaire sur le $A_{\mathfrak p}$-module $N_{\mathfrak p}$ s'écrit
$u=v/f$, où $f\in A\setmin\mathfrak p$, et où $v$ est une forme linéaire sur le
$A$-module $N$. Il est clair que \emph{c)} entraîne \emph{b)}, et il reste donc à voir que \emph{a)}
entraîne \emph{d)}. Or, comme $N$ est de présentation finie,
$(\operatorname{Hom}_A(N,M))_{\mathfrak p}$ s'identifie canoniquement à
$\operatorname{Hom}_{A_{\mathfrak p}}(N_{\mathfrak p},M_{\mathfrak p})$
(Bourbaki, \emph{Alg. comm.}, chap.~II, §~2, n\textsuperscript{o}~7,
prop.~19). Si $w_{\mathfrak p}$ est inverse à gauche de $u_{\mathfrak p}$, il
existe donc un homomorphisme $w:N\to M$ et un élément $f\in A\setmin\mathfrak p$
tels que $w=w_{\mathfrak p}\otimes(1/f)1_{A_{\mathfrak p}}$. La relation
$w_{\mathfrak p}\circ u_{\mathfrak p}=1_{M_{\mathfrak p}}$ s'écrit donc aussi
$(w\circ u)\otimes1_{A_{\mathfrak p}}=f.1_{M_{\mathfrak p}}$. Mais comme $M$
est un $A$-module de type fini, il existe $g\in A\setmin\mathfrak p$ tel que
l'endomorphisme $g((w\circ u)\setmin f.1_M)$ s'annule en tous les générateurs de $M$,
donc soit nul. Si on pose $h=fg$, et
$u_h=u\otimes1_{A_h}$, $w_h=w\otimes1_{A_h}$, on a donc
$w_h(u_h(z))=(f/1)z$ pour tout $z\in M_h$~; mais comme $h/1$ est inversible
dans $A_h$, il en est de même de $f/1=f'$, et $f'^{-1}w_h$ est par suite
inverse à gauche de $u_h$, ce qui achève la démonstration.

\begin{proposition}[19.1.13]
\label{0.19.1.13-fr}
Supposons vérifiées les hypothèses de \textup{(19.1.9)} et supposons en outre que
$(\mathfrak J_\lambda)$ soit une suite décroissante $(\mathfrak J_n)$ et que
$M$ soit séparé et complet. Alors les conditions \textup{a)} et \textup{b)} de \textup{(19.1.9)} sont
aussi équivalentes à~:
\begin{enumerate}
  \item[\textup{c)}] $u$ est inversible à gauche.
\end{enumerate}
\end{proposition}

On sait déjà que \emph{c)} entraîne \emph{a)} (19.1.6). Inversement, si \emph{a)} est vérifiée, on
sait (avec les notations de (19.1.8)) que $M_n$ est facteur direct du
$(A/\mathfrak J_n)$-module projectif $N_n$, donc $P_n$ est aussi isomorphe à
un facteur direct de $N_n$, et par suite est projectif~; il suffit donc
d'appliquer (19.1.8).

Notons enfin la proposition suivante~:

\begin{proposition}[19.1.14]
\label{0.19.1.14-fr}
Soient $A$ un anneau, $M$ un $A$-module de type fini, $N$ un $A$-module
projectif, $u:M\to N$ un $A$-homomorphisme.
\begin{enumerate}
  \item[\textup{(i)}] Pour que $u$ soit inversible à gauche, il faut et il suffit que,
  pour tout idéal maximal $\mathfrak m$ de $A$,
  $u_{\mathfrak m}:M_{\mathfrak m}\to N_{\mathfrak m}$ soit inversible à
  gauche.
  \item[\textup{(ii)}] Soit $A'$ une $A$-algèbre qui soit un $A$-module fidèlement
  plat. Pour que $u$ soit inversible à gauche, il faut et il suffit que
  $u\otimes1:M\otimes_AA'\to N\otimes_AA'$ soit inversible à gauche.
\end{enumerate}
\end{proposition}

Comme dans (19.1.12), on peut se borner au cas où $N$ est libre de type fini~;
dire que $u$ est inversible à gauche signifie alors que $u$ est injectif et
que le module quotient $P=N/u(M)$ est \emph{projectif}, car $u(M)$ sera alors
facteur direct de $N$. Notons en outre que puisque $M$ est de type fini, $P$
est de \emph{présentation finie}. Cela étant~:

(i) La condition est évidemment nécessaire. Inversement, si elle est
satisfaite, on sait que $u$ est injectif (Bourbaki, \emph{Alg. comm.},
chap.~II, §~3, n\textsuperscript{o}~3, th.~1) et comme
$P_{\mathfrak m}=N_{\mathfrak m}/u_{\mathfrak m}(M_{\mathfrak m})$ est
projectif pour tout $\mathfrak m$, on sait que cela entraîne que $P$ est
projectif (\emph{loc. cit.}, §~5, n\textsuperscript{o}~2, th.~1).

(ii) Ici encore, la condition est trivialement nécessaire. Inversement, si
elle est remplie, on sait que $u$ est injectif ($\mathbf{0}_{\mathrm I}$,
6.4.1) et comme $P\otimes_AA'=\operatorname{Coker}(u\otimes1)$ est projectif,
donc plat, on en déduit que $P$ est un $A$-module plat
($\mathbf{0}_{\mathrm I}$, 6.6.3), donc projectif puisqu'il est de présentation
finie (Bourbaki, \emph{Alg. comm.}, chap.~II, §~5,
n\textsuperscript{o}~2, cor.~2 du th.~1).

\oldpage[0\textsubscript{IV-1}]{78}

\begin{remark}[19.1.15]
\label{0.19.1.15-fr}
La notion «~duale~» de celle d'homomorphisme formellement inversible à gauche
est celle d'homomorphisme \emph{formellement inversible à droite}~; un tel
$A$-homomorphisme continu $u:M\to N$ vérifie, par définition, la condition
suivante~: pour tout sous-module ouvert $W'$ de $N$, il existe un sous-module
ouvert $V'\subset u^{-1}(W')$ de $M$, un sous-module ouvert $W''\subset W'$ de
$N$ et un homomorphisme $h:N/W''\to M/V'$ tels que l'homomorphisme canonique
$N/W''\to N/W'$ se factorise en
\[
  N/W''\xrightarrow{h}M/V'\xrightarrow{u'}N/W'
\]
où $u'$ est déduit de $u$ par passage aux quotients. Cela entraîne que $u$ est
un épimorphisme formel~; s'il existe un $\widehat A$-homomorphisme continu
$r:\widehat N\to\widehat M$ tel que $\widehat u\circ r=1_{\widehat N}$, on
vérifie aussitôt que $u$ est formellement inversible à droite. Si les
conditions de (19.1.7) sont remplies, pour que $u$ soit formellement inversible
à droite, il faut et il suffit que les $u_\lambda$ soient inversibles à droite
(c'est-à-dire que le noyau de $u_\lambda$ est facteur direct de
$M/\mathfrak J_\lambda M$ et que $u_\lambda$ est un isomorphisme d'un
supplémentaire de $\operatorname{Ker}(u_\lambda)$ sur
$N/\mathfrak J_\lambda N$). Nous laissons au lecteur le soin d'énoncer et de
démontrer les propositions correspondant à (19.1.8) et (19.1.9) (dans
l'analogue de (19.1.8), il faut supposer $M$ séparé et complet et que $M_n$ est
un $(A/\mathfrak J_n)$-module projectif~; dans l'analogue de (19.1.9), il faut
supprimer l'hypothèse sur les $N/\mathfrak J_\lambda N$, mais supposer par
contre que, pour tout $\lambda$, $M/\mathfrak J_\lambda M$ est un
$(A/\mathfrak J_\lambda)$-module projectif).
\end{remark}

\subsection{Modules formellement projectifs.}
\label{subsection:0.19.2-fr}

\begin{definition}[19.2.1]
\label{0.19.2.1-fr}
Soient $A$ un anneau topologique. On dit qu'un $A$-module topologique $M$ est
formellement projectif s'il vérifie la condition suivante~: pour tout
idéal ouvert $\mathfrak J$ de $A$, tout couple de $(A/\mathfrak J)$-modules
(discrets) $P$, $Q$, tout $A$-homomorphisme surjectif $u:P\to Q$ et tout
$A$-homomorphisme continu $v:M\to Q$, il existe un $A$-homomorphisme continu
$w:M\to P$ tel que $v=u\circ w$.
\end{definition}

\begin{env}[19.2.2]
\label{0.19.2.2-fr}
Pour vérifier la condition de (19.2.1), on peut évidemment se borner (en
remplaçant $Q$ par $v(M)$ et $P$ par $u^{-1}(v(M))$) au cas où $v$ est
lui-même surjectif~; alors $Q$ est isomorphe à $M/V$, où $V$ est un
sous-module ouvert de $M$ tel que $\mathfrak JM\subset V$~; la condition de
(19.2.1) équivaut alors à dire que pour tout $(A/\mathfrak J)$-module discret
$P$ et tout $A$-homomorphisme surjectif $u:P\to M/V$, il existe un sous-module
ouvert $V'\subset V$ dans $M$ et un $A$-homomorphisme $w:M/V'\to P$, tels que
l'homomorphisme canonique $M/V'\to M/V$ se factorise en
$M/V'\xrightarrow{w}P\xrightarrow{u}M/V$.
On notera qu'il suffit de vérifier cette condition lorsque $\mathfrak J$
parcourt un système fondamental de voisinages de 0 dans $A$ formé d'idéaux.
\end{env}

\begin{env}[19.2.3]
\label{0.19.2.3-fr}
Supposons qu'il existe un système fondamental $(\mathfrak J_\lambda)$ d'idéaux
ouverts dans $A$ et un système fondamental $(V_\lambda)$ de sous-modules
ouverts de $M$, ayant \emph{même} ensemble d'indices que
$(\mathfrak J_\lambda)$ et tel que, pour tout $\lambda$, $M/V_\lambda$ soit un
$(A/\mathfrak J_\lambda)$-module \emph{projectif}. Alors $M$ est formellement
projectif~: il suffit en effet, avec les notations de (19.2.2), de prendre
$\lambda$ tel que $\mathfrak J_\lambda\subset\mathfrak J$ et
$V_\lambda\subset V$~; comme $P$ est aussi un
$(A/\mathfrak J_\lambda)$-module, la factorisation de l'homomorphisme canonique
$M/V_\lambda\to M/V$ en $M/V_\lambda\to P\to M/V$ résulte de ce qu'il s'agit
d'un $(A/\mathfrak J_\lambda)$-homomorphisme et de ce que $M/V_\lambda$ est
supposé être un $(A/\mathfrak J_\lambda)$-module projectif.

Lorsque la condition plus stricte de ce numéro est satisfaite, on dit que $M$
est \emph{strictement formellement projectif}.
\end{env}

\begin{proposition}[19.2.4]
\label{0.19.2.4-fr}
Soient $A$ un anneau topologique, $M$ un $A$-module topologique dont la
topologie est déduite de celle de $A$. Pour que $M$ soit formellement
projectif, il faut et il suffit qu'il soit strictement formellement projectif.
\end{proposition}

\oldpage[0\textsubscript{IV-1}]{79}

\begin{proof}
On vient de voir que la condition est suffisante. Inversement, supposons $M$
formellement projectif et soit $(\mathfrak J_\lambda)$ un système fondamental
d'idéaux ouverts dans $A$. Pour tout $\lambda$, soient $P$ un
$(A/\mathfrak J_\lambda)$-module libre et
$u:P\to M/\mathfrak J_\lambda M$ un
$(A/\mathfrak J_\lambda)$-homomorphisme surjectif. Il existe donc un
$\mathfrak J_\mu\subset\mathfrak J_\lambda$ tel que l'homomorphisme canonique
$M/\mathfrak J_\mu M\to M/\mathfrak J_\lambda M$ se factorise en
$M/\mathfrak J_\mu M\xrightarrow{w}P\xrightarrow{u}M/\mathfrak J_\lambda M$~;
mais comme
$w(\mathfrak J_\lambda(M/\mathfrak J_\mu M))\subset\mathfrak J_\lambda P=0$,
$w$ se factorise en $M/\mathfrak J_\mu M\to M/\mathfrak J_\lambda M\xrightarrow{v}P$, et il est
clair que $u\circ v$ est l'identité dans $M/\mathfrak J_\lambda M$, ce qui
prouve que ce $(A/\mathfrak J_\lambda)$-module est projectif.
\end{proof}

\begin{proposition}[19.2.5]
\label{0.19.2.5-fr}
Soient $A$ un anneau topologique, $M$ un $A$-module topologique.
\begin{enumerate}
  \item[\textup{(i)}] Pour que $M$ soit formellement projectif \textup{(resp.} strictement
  formellement projectif\textup{)}, il faut et il suffit que le $\hat A$-module
  topologique $\hat M$ le soit.
  \item[\textup{(ii)}] Soit $A'$ une $A$-algèbre topologique. Si $M$ est formellement
  projectif \textup{(resp.} strictement formellement projectif\textup{)}, alors
  $M\otimes_A A'$ (muni de la topologie produit tensoriel) est un $A'$-module
  topologique formellement projectif \textup{(resp.} strictement formellement
  projectif\textup{)}.
\end{enumerate}
\end{proposition}

\begin{proof}
\begin{enumerate}
  \item[(i)] Il suffit de remarquer que lorsque $\mathfrak J$ (resp. $V$)
  parcourt l'ensemble des idéaux ouverts de $A$ (resp. l'ensemble des
  sous-modules ouverts de $M$), le séparé complété $\hat{\mathfrak J}$
  (resp. $\hat V$) parcourt l'ensemble des idéaux ouverts de $\hat A$
  (resp. l'ensemble des sous-modules ouverts de $\hat M$), et
  $\hat A/\hat{\mathfrak J}=A/\mathfrak J$ (resp.
  $\hat M/\hat V=M/V$) à un isomorphisme canonique près~; comme la notion
  de module formellement projectif (resp. strictement formellement projectif)
  ne fait intervenir que les $A/\mathfrak J$ et les $M/V$, on en déduit
  aussitôt (i).
  \item[(ii)] Supposons d'abord $M$ formellement projectif et posons
  $M'=M\otimes_A A'$~; soient $\mathfrak J'$ un idéal ouvert de $A'$, $P'$,
  $Q'$ deux $(A'/\mathfrak J')$-modules discrets, $u':P'\to Q'$ un
  $A'$-homomorphisme surjectif, $v':M'\to Q'$ un $A'$-homomorphisme continu.
  Il y a un idéal ouvert $\mathfrak J$ de $A$ tel que
  $\mathfrak JA'\subset\mathfrak J'$, donc $P'$ et $Q'$ sont aussi des
  $(A/\mathfrak J)$-modules discrets. Si l'on considère le
  $A$-homomorphisme composé $v:M\xrightarrow{j}M'\xrightarrow{v}Q'$ qui est
  continu, l'hypothèse entraîne qu'il existe un $A$-homomorphisme continu
  $w:M\to P'$ tel que $v=u'\circ w$~; mais comme $P'$ est un $A'$-module
  topologique, $w$ se factorise en
  $M\xrightarrow{j}M'\xrightarrow{w'}P'$, où $w'$ est un $A'$-homomorphisme
  continu, et comme $v'\circ j=(u'\circ w')\circ j$, on en conclut que
  $v'=u'\circ w'$.

  Supposons ensuite que $M$ soit strictement formellement projectif~; soit
  $(\mathfrak J_\lambda)$ (resp. $(V_\lambda)$) un système fondamental
  d'idéaux ouverts (resp. de sous-modules ouverts) dans $A$ (resp. $M$) tel
  que $(M/V_\lambda)$ soit un $(A/\mathfrak J_\lambda)$-module projectif,
  et soit $(\mathfrak J'_\mu)$ un système fondamental d'idéaux ouverts dans
  $A'$. On a un système fondamental de voisinages de $0$ dans $M'$ en prenant
  les sous-modules
  $\operatorname{Im}(M\otimes_A\mathfrak J'_\mu)
    +\operatorname{Im}(V_\lambda\otimes_A A')=W_{\lambda\mu}$,
  où $\lambda$ et $\mu$ sont tels que
  $\mathfrak J_\lambda A'\subset\mathfrak J'_\mu$. Comme
  $(M\otimes_A A')/W_{\lambda\mu}
    =(M/V_\lambda)\otimes_{A/\mathfrak J_\lambda}
      (A'/\mathfrak J'_\mu)$
  et que $(M/V_\lambda)$ est un
  $(A/\mathfrak J_\lambda)$-module projectif, $M'/W_{\lambda\mu}$
  est un $(A'/\mathfrak J'_\mu)$-module projectif.
\end{enumerate}
\end{proof}

\subsection{Algèbres formellement lisses.}
\label{subsection:0.19.3-fr}

\begin{definition}[19.3.1]
\label{0.19.3.1-fr}
Soient $A$ un anneau topologique, $B$ une $A$-algèbre topologique. On dit que
$B$ est une $A$-algèbre formellement lisse si, pour toute
$A$-algèbre topologique \textup{discrète} $C$
\oldpage[0\textsubscript{IV-1}]{80}
et tout idéal \textup{nilpotent} $\mathfrak J$ de
$C$, tout $A$-homomorphisme continu $u:B\to C/\mathfrak J$ se factorise en
$B\xrightarrow{v}C\xrightarrow{\varphi}C/\mathfrak J$
où $v$ est un homomorphisme continu et $\varphi$ l'homomorphisme canonique.
\end{definition}

La définition (19.3.1) revient à dire que la propriété suivante a lieu~:
\begin{enumerate}
  \item[(P)] Pour tout idéal \emph{ouvert} $\mathfrak K$ de la $A$-algèbre $B$ et
  tout $A$-homomorphisme $u':B/\mathfrak K\to C/\mathfrak J$, il y a un
  idéal ouvert $\mathfrak K'\subset\mathfrak K$ de $B$ tel que
  l'homomorphisme
  $B\to B/\mathfrak K\xrightarrow{u'}C/\mathfrak J$ se factorise en
  $B\to B/\mathfrak K'\xrightarrow{v'}C\to C/\mathfrak J$,
  où $v'$ est un $A$-homomorphisme.
\end{enumerate}

En effet, si $u:B\to C/\mathfrak J$ est un $A$-homomorphisme continu, il a
pour noyau $\mathfrak K$ un idéal ouvert de $B$, donc $u$ se factorise en
$B\to B/\mathfrak K\xrightarrow{u'}C/\mathfrak J$, et si (P) est
vérifiée, il suffit de l'appliquer à $u'$ et de prendre pour $v$ la composée
$B\to B/\mathfrak K'\xrightarrow{v'}C$ pour satisfaire aux conditions de
(19.3.1). Inversement, supposons que $B$ soit une $A$-algèbre formellement
lisse~; donnons-nous un idéal ouvert $\mathfrak K$ de $B$ et un
$A$-homomorphisme $u':B/\mathfrak K\to C/\mathfrak J$, et appliquons la
définition (19.3.1) à
$u:B\to B/\mathfrak K\xrightarrow{u'}C/\mathfrak J$~; si
$v:B\to C$ est un $A$-homomorphisme continu tel que $u$ se factorise en
$B\xrightarrow{v}C\to C/\mathfrak J$, l'idéal
$\mathfrak K'=\operatorname{Ker}(v)\cap\mathfrak K$ de $B$ est ouvert et
contenu dans $\mathfrak K$~; par suite $v$ se factorise en
$B\to B/\mathfrak K'\xrightarrow{v'}C$, et $v'$ vérifie bien la condition
(P).

\begin{proposition}[19.3.2]
\label{0.19.3.2-fr}
Soient $A$ un anneau discret, $V$ un $A$-module projectif~; l'algèbre
symétrique $B=\mathbf S_A^\bullet(V)$, munie de la topologie discrète, est une
$A$-algèbre formellement lisse.
\end{proposition}

\begin{proof}
En effet, les notations $C$ et $\mathfrak J$ ayant le même sens que dans
(19.3.1), soit $u:B\to C/\mathfrak J$ un homomorphisme de $A$-algèbres, qui
par restriction à $V=\mathbf S_A^1(V)$, donne un homomorphisme de $A$-modules
$u_1:V\to C/\mathfrak J$~; comme $V$ est projectif, $u_1$ se factorise en
$V\xrightarrow{v_1}C\xrightarrow{\varphi}C/\mathfrak J$, et $v_1$ se
prolonge en un homomorphisme de $A$-algèbres
$v:\mathbf S_A^\bullet(V)\to C$, tel que le composé
$\varphi\circ v$ coïncide avec $u$ dans $V$, donc
$\varphi\circ v=u$, ce qui démontre la proposition.
\end{proof}

\begin{corollary}[19.3.3]
\label{0.19.3.3-fr}
Si $A$ est un anneau discret, toute algèbre de polynômes
$B=A[T_\alpha]_{\alpha\in I}$, munie de la topologie discrète, est une
$A$-algèbre formellement lisse.
\end{corollary}

\begin{proposition}[19.3.4]
\label{0.19.3.4-fr}
Soit $A$ un anneau topologique, et soit
$B=A[[T_\alpha]]_{\alpha\in I}$ un anneau de séries formelles
$\sum_{\lambda\in\mathbb{N}^{(I)}}c_\lambda T^\lambda$ (algèbre large sur
$A$ du monoïde $\mathbb{N}^{(I)}$, identifié en tant que $A$-module au produit
$A^{\mathbb{N}^{(I)}}$). Si l'on munit $B$ de la topologie produit, $B$ est une
$A$-algèbre formellement lisse.
\end{proposition}

Soit $(\mathfrak L_\mu)_{\mu\in M}$ un système fondamental d'idéaux ouverts
dans $A$. Pour toute partie finie $H$ de $I$, tout $\mu\in M$ et tout entier
$n$, désignons par $\mathfrak K_{H,\mu,n}$ le voisinage de $0$ dans $B$ formé
des $(c_\lambda)$ tels que pour tout
$\lambda=(\lambda_\alpha)_{\alpha\in I}$ tel que
$\lambda_\alpha=0$ pour $\alpha\notin H$ et
$\sum_{\alpha\in H}\lambda_\alpha\leq n$, on ait
$c_\lambda\in\mathfrak L_\mu$. On vérifie immédiatement que les
$\mathfrak K_{H,\mu,n}$ forment un système fondamental de voisinages de $0$
dans $B$ et sont des idéaux de $B$, donc la topologie produit est compatible
avec la structure de $A$-algèbre de $B$.

Notons d'abord, avec les mêmes notations, le
\begin{lemma}[19.3.4.1]
\label{0.19.3.4.1-fr}
Soit $E$ une $A$-algèbre discrète.
\begin{enumerate}
  \item[\textup{(i)}] Si $f:B\to E$ est un $A$-homomorphisme continu, il existe une
  partie finie $H$ de $I$ telle que $f(T_\alpha)=0$ pour $\alpha\notin H$,
  et $f(T_\alpha)$ est nilpotent dans $E$ pour tout $\alpha\in H$.
  \oldpage[0\textsubscript{IV-1}]{81}
  \item[\textup{(ii)}] Inversement, soient $H$ une partie finie de $I$,
  $(z_\alpha)_{\alpha\in H}$ une famille d'éléments nilpotents de $E$.
  Il existe un $A$-homomorphisme continu $g:B\to E$ et un seul tel que
  $g(T_\alpha)=z_\alpha$ pour $\alpha\in H$ et
  $g(T_\alpha)=0$ pour $\alpha\notin H$.
\end{enumerate}
\end{lemma}

\begin{proof}
(i) résulte aussitôt de ce que $f^{-1}(0)$ est un voisinage de $0$ dans le
$A$-module produit $A^{\mathbb{N}^{(I)}}$, d'où $f(T^\lambda)=0$ sauf pour un
nombre fini de valeurs de $\lambda\in\mathbb{N}^{(I)}$. Pour prouver (ii), il
suffit de remarquer que l'anneau de polynômes
$B'=A[T_\alpha]_{\alpha\in I}$ est dense dans $B$~; l'existence et
l'unicité de la restriction $g|B'$ sont triviales et sa continuité résulte de
l'hypothèse que les $f(T_\alpha)$ pour $\alpha\in H$ sont nilpotents,
car si $(f(T_\alpha))^n=0$ pour tout $\alpha\in H$, on a
$g(T^\lambda)=0$ pour tout $\lambda=(\lambda_\alpha)_{\alpha\in I}$ de
support fini sauf ceux tels que $\lambda_\alpha=0$ pour
$\alpha\notin H$ et $\lambda_\alpha\leq n$ pour
$\alpha\in H$, c'est-à-dire sauf pour un nombre fini de valeurs de
$\lambda$.
\end{proof}

\begin{proof}[Démonstration de la proposition~{\rm(19.3.4)}]
Ce lemme étant établi, et les notations $C$ et $\mathfrak J$ ayant le même
sens que dans (19.3.1), on a $u(T_\alpha)=0$ sauf pour
$\alpha\in H$, $H$ étant une partie finie de $I$, et les
$z_\alpha=u(T_\alpha)$ pour $\alpha\in H$ sont nilpotents dans
$C/\mathfrak J$~; comme $\mathfrak J$ est nilpotent, il existe une famille
$(x_\alpha)_{\alpha\in H}$ d'éléments nilpotents de $C$ dont les images
canoniques dans $C/\mathfrak J$ sont les $z_\alpha$~; si $v$ est le
$A$-homomorphisme continu de $B$ dans $C$ tel que
$v(T_\alpha)=0$ pour $\alpha\notin H$, $v(T_\alpha)=x_\alpha$ pour
$\alpha\in H$, il est clair que $u$ se factorise en
$B\xrightarrow{v}C\to C/\mathfrak J$.
\end{proof}

\begin{proposition}[19.3.5]
\label{0.19.3.5-fr}
Soit $A$ un anneau topologique.
\begin{enumerate}
  \item[(i)] $A$ est une $A$-algèbre formellement lisse.
  \item[(ii)] Si $B$ est une $A$-algèbre formellement lisse et $C$ une
  $B$-algèbre formellement lisse, alors $C$ est une $A$-algèbre formellement
  lisse.
  \item[(iii)] Soient $B$ une $A$-algèbre formellement lisse, $A'$ une
  $A$-algèbre topologique~; alors la $A'$-algèbre topologique
  $B\otimes_A A'$ ($\mathbf 0_{\rm I}$, 7.7.5 et 7.7.6) est formellement
  lisse.
  \item[(iv)] Soient $B$ une $A$-algèbre topologique, $S$ (resp. $T$) une
  partie multiplicative de $A$ (resp. $B$) telle que l'image canonique de $S$
  dans $B$ soit contenue dans $T$. Si $B$ est une $A$-algèbre formellement
  lisse, alors $T^{-1}B$ est une $S^{-1}A$-algèbre formellement lisse.
  \item[(v)] Soient $B_i$ ($1\leq i\leq n$) des $A$-algèbres topologiques.
  Pour que $\prod_{i=1}^n B_i$ soit une $A$-algèbre formellement lisse, il
  faut et il suffit que chacune des $B_i$ le soit.
\end{enumerate}
\end{proposition}

